% ---------------------------------------------------------------------------
% File trees for A1, A2 and A5, plus the shared ExampleCode folder.
%
% Needs in the preamble:   \usepackage{dirtree}
% Then in each section:    \ExampleCodeTree   and   \AOneTree / \ATwoTree / \AFiveTree
%
% dirtree syntax: ".<level> <name>." with \DTcomment{...} for the note.
% ---------------------------------------------------------------------------

% ---- ExampleCode: shared by every assignment, built from ../ExampleCode ----
\newcommand{\ExampleCodeTree}{%
\dirtree{%
.1 ExampleCode/\DTcomment{task-supplied library and has been modified to be used for A1, A2 and A5}.
.2 CRender.h\DTcomment{the only class allowed to talk to raylib; draws loops, discs, lines, trails}.
.2 CRender.cpp\DTcomment{raylib calls live here; added DrawTrail and silenced raylib's log}.
.2 CLoopReader.h\DTcomment{reads a .map file into a name, a start pose and a closed loop of vertices}.
.2 CLoopReader.cpp\DTcomment{the parser}.
.2 SimpleWalls.map\DTcomment{the L-shaped room the wall follower drives round}.
.2 SimpleLine.map\DTcomment{the closed loop on the floor the line follower tracks}.
.2 TestRender.cpp\DTcomment{staff demo of CRender, not used by the simulator}.
.2 TestLoopReader.cpp\DTcomment{staff demo of CLoopReader, not used by the simulator}.
.2 readme\DTcomment{staff notes on installing raylib and building the demos}.
}}

% ---- A1: one wall follower ----
\newcommand{\AOneTree}{%
\dirtree{%
.1 A1/.
.2 CMakeLists.txt\DTcomment{builds simulator from src/ and ../ExampleCode, output to bin/}.
.2 include/.
.3 CSimulator.h\DTcomment{owns the map, renderer and clock; steps and draws every robot}.
.3 CRobot.h\DTcomment{base robot: pose, two motors, kinematics, trail, solid-wall collision handling}.
.3 CWFRobot.h\DTcomment{wall follower: two range sensors and the two-term controller}.
.3 CSensor.h\DTcomment{range sensor: casts a ray at an angle and reports the distance}.
.3 CMotor.h\DTcomment{one wheel; holds a speed clamped to its limits}.
.3 Geometry.h\DTcomment{Vec2D, CPose, ray--wall intersection, nearest point and distance to a wall}.
.3 Helper.h\DTcomment{ClampInt}.
.2 src/.
.3 main.cpp\DTcomment{loads SimpleWalls.map, makes one CWFRobot, runs the simulator}.
.3 CSimulator.cpp.
.3 CRobot.cpp\DTcomment{Drive() with midpoint integration; HandleCollision() pushes the robot back inside and counts contacts}.
.3 CWFRobot.cpp\DTcomment{Update(): distance and heading errors to a turn rate}.
.3 CSensor.cpp.
.3 CMotor.cpp.
.3 Geometry.cpp.
.3 Helper.cpp.
}}

% ---- A2: wall follower and line follower together ----
\newcommand{\ATwoTree}{%
\dirtree{%
.1 A2/.
.2 CMakeLists.txt\DTcomment{as A1, with the new sensor and robot sources added}.
.2 include/.
.3 CSimulator.h\DTcomment{now borrows any number of maps and robots; only draws and steps}.
.3 CRobot.h\DTcomment{base robot; given the map it senses and the room it can hit, plus its colours}.
.3 CWFRobot.h\DTcomment{wall follower, unchanged controller; senses and collides with the room}.
.3 CLFRobot.h\DTcomment{line follower: two line sensors on the rim and a two-bit controller}.
.3 CSensor.h\DTcomment{sensor base: owns the mounting and the robot-to-world transform}.
.3 CWallSensor.h\DTcomment{range sensor, was A1's CSensor; returns a distance}.
.3 CLineSensor.h\DTcomment{line sensor at an offset; returns whether it is over the line}.
.3 CMotor.h\DTcomment{unchanged}.
.3 Geometry.h\DTcomment{same maths as A1; used by both sensors and the collision handling}.
.3 Helper.h\DTcomment{Clamp, overloaded for int and float}.
.2 src/.
.3 main.cpp\DTcomment{loads both maps, makes one robot of each type, lends them to the simulator}.
.3 CSimulator.cpp.
.3 CRobot.cpp.
.3 CWFRobot.cpp.
.3 CLFRobot.cpp\DTcomment{Update(): sensor bits to a turn rate}.
.3 CSensor.cpp.
.3 CWallSensor.cpp.
.3 CLineSensor.cpp.
.3 CMotor.cpp.
.3 Geometry.cpp.
.3 Helper.cpp.
}}

% ---- A5: twenty of each, with noise ----
\newcommand{\AFiveTree}{%
\dirtree{%
.1 A5/\DTcomment{a copy of A2; only the files noted below differ}.
.2 CMakeLists.txt.
.2 include/.
.3 CSimulator.h.
.3 CRobot.h\DTcomment{adds RandomOffset() and the three noise constants}.
.3 CWFRobot.h.
.3 CLFRobot.h.
.3 CSensor.h.
.3 CWallSensor.h.
.3 CLineSensor.h.
.3 CMotor.h.
.3 Geometry.h.
.3 Helper.h.
.2 src/.
.3 main.cpp\DTcomment{seeds rand(), makes twenty of each robot in reserved vectors, one colour each; two flags pick which types run}.
.3 CSimulator.cpp.
.3 CRobot.cpp\DTcomment{noise on the start pose in the constructor and on each wheel in Drive()}.
.3 CWFRobot.cpp.
.3 CLFRobot.cpp.
.3 CSensor.cpp.
.3 CWallSensor.cpp.
.3 CLineSensor.cpp.
.3 CMotor.cpp.
.3 Geometry.cpp.
.3 Helper.cpp.
}}
